\documentclass{ximera}
%verbeterd 12/11/26

%   \author{}
    \title{Bibliografie}
\begin{document}
\begin{abstract}
    Lijst van belangrijkste gebruikte bronnen.
\end{abstract}
\maketitle
    

\begin{itemize}
    \item Thomson Learning International Edition, \textit{Physics for scientists and engineers with Modern Physics}, Seventh edition, R. A. Serway and J.W. Jewett.
 
    \item Interactie 5.1, Die Keure
\end{itemize}
 
 
 

\section*{Bronnen van afbeeldingen}
Figuren zonder bronvermelding zijn eigen tekeningen (TikZ) van de auteur. Foto's komen van Wikimedia Commons en worden gebruikt onder de vermelde licentie: publiek domein, of een Creative Commons-licentie (CC BY of CC BY-SA, zie \url{https://creativecommons.org/licenses/}).

\subsection*{Het magnetisch veld}
\begin{itemize}
        \item Figuur~\verwijs[Magnetiet trekt spijkers aan]{fig:magnetiet}{magnetisme_permanente_magneten}: F. A. Annet (1922), publiek domein, Wikimedia Commons, \url{https://commons.wikimedia.org/wiki/File:Lodestone_attracting_nails.png}
        \item Figuur~\verwijs[IJzervijlsel rond een staafmagneet]{fig:ijzervijlsel_staafmagneet}{magnetisme_veld}: MikeRun en Vera Wurmsdobler, CC BY-SA 4.0, Wikimedia Commons, \url{https://commons.wikimedia.org/wiki/File:Bar-magnet-iron-filings_max.jpg}
        \item Figuur~\verwijs[IJzervijlsel tussen twee staafmagneten (a)]{fig:ijzervijlsel_twee_magneten}{magnetisme_veld}: A. W. Duff (1908), publiek domein, Wikimedia Commons, \url{https://commons.wikimedia.org/wiki/File:Magnetic_field_of_bar_magnets_attracting.png}
        \item Figuur~\verwijs[IJzervijlsel tussen twee staafmagneten (b)]{fig:ijzervijlsel_twee_magneten}{magnetisme_veld}: A. W. Duff (1908), publiek domein, Wikimedia Commons, \url{https://commons.wikimedia.org/wiki/File:Magnetic_field_of_bar_magnets_repelling.png}
        \item Figuur~\verwijs[Schooltoestel voor de proef van Oersted]{fig:proef_oersted}{magnetisme_rechte_geleider}: A. Privat-Deschanel (19de eeuw), publiek domein, Wikimedia Commons, \url{https://commons.wikimedia.org/wiki/File:Oersted_experiment.png}
        \item Figuur~\verwijs[IJzervijlsel rond een stroomvoerende draad]{fig:ijzervijlsel_rechte_geleider}{magnetisme_rechte_geleider}: Maciej J. Mrowinski, CC BY-SA 4.0, Wikimedia Commons, \url{https://commons.wikimedia.org/wiki/File:Magnetic_field_around_wire.jpg}
        \item Figuur~\verwijs[IJzervijlsel rond een stroomvoerende spoel]{fig:ijzervijlsel_spoel}{magnetisme_spoel}: Maciej J. Mrowinski, CC BY-SA 4.0, Wikimedia Commons, \url{https://commons.wikimedia.org/wiki/File:Magnetic_field_around_solenoid.jpg}
        \item Figuur~\verwijs[Hijsmagneet]{fig:hijsmagneet}{magnetisme_elektromagneet}: C. M. Jansky (1914), publiek domein, Wikimedia Commons, \url{https://commons.wikimedia.org/wiki/File:Industrial_lifting_magnet.jpg}
        \item Figuur~\verwijs[MRI-scanner]{fig:mri_scanner}{magnetisme_elektromagneet}: Jan Ainali, CC BY 3.0, Wikimedia Commons, \url{https://commons.wikimedia.org/wiki/File:MRI-Philips.JPG}
\end{itemize}

\end{document}